\documentclass[aspectratio=169]{beamer}
\usepackage[utf8]{inputenc}
\usepackage[T1]{fontenc}
\usepackage[french]{babel}
\usepackage{amsmath}
\usepackage{amssymb}
\usepackage{graphicx}
\usepackage{pgfpages}
\usepackage{tikz}
\usetikzlibrary{arrows.meta, positioning, decorations.pathreplacing, calc, shapes.geometric, fit, backgrounds}
\usepackage[backend=biber, style=alphabetic, url=true]{biblatex}
\addbibresource{../../biblio.bib}

\graphicspath{{images/}}
\usetheme{Boadilla}
\usecolortheme{seahorse}
% \setbeameroption{show notes on second screen=right}

\title{Leçon 6 : Vérification formelle à l'ère de l'IA}
\subtitle{Promesses, retournements, questions ouvertes}
\author{Pablo Donato}
\institute{Logique et Fondements de l'Informatique}
\date{Année 2026}

\begin{document}

% ─────────────────────────────────────────────────────────────
% FRAME 1 — TITRE
% ─────────────────────────────────────────────────────────────
\begin{frame}
  \titlepage
  \note{
    \textbf{Cadre}
    \begin{itemize}
      \item Dernière séance ; pas de nouveau concept technique.
      \item On prend du recul philosophique sur les 6 semaines.
    \end{itemize}
    \textbf{Fil dialectique}
    \begin{itemize}
      \item Acte I --- la promesse (maths, sûreté, IA safety).
      \item Acte II --- le retournement 2026 (trois fractures).
      \item Acte III --- questions ouvertes, discussion.
    \end{itemize}
  }
\end{frame}

% ─────────────────────────────────────────────────────────────
% ACTE I — LA PROMESSE
% ─────────────────────────────────────────────────────────────
\section{Acte I — La promesse}

% ─────────────────────────────────────────────────────────────
% FRAME 2 — DE LA VÉRIFICATION CLASSIQUE AU RÊVE LEIBNIZIEN
% ─────────────────────────────────────────────────────────────
\begin{frame}{Du rêve de Leibniz à la vérification formelle}
  \begin{columns}
    \begin{column}{0.55\textwidth}
      \textbf{Leibniz, XVII\textsuperscript{e} :} \emph{characteristica universalis} + \emph{calculus ratiocinator}
      \medskip

      \textbf{Réalisations :}
      \begin{itemize}
        \item seL4~\cite{klein2009sel4} — noyau OS prouvé
        \item CompCert~\cite{leroy2009compcert} — compilateur C certifié
        \item 4 Couleurs~\cite{gonthier2008fourcolor}, Kepler~\cite{hales2017flyspeck}
      \end{itemize}
      \medskip
      \small\emph{Mais : 20~années-personnes pour seL4, 6 pour CompCert.}
    \end{column}
    \begin{column}{0.42\textwidth}
      \begin{block}{Le pari}
        Garanties \emph{mathématiques} là où l'ingénierie classique n'en offre pas.
      \end{block}
      \medskip
      \small
      Benzmüller~\cite{benzmullerUniversalReasoningRational2017} :
      logique d'ordre supérieur comme \emph{meta-logique universelle}.
    \end{column}
  \end{columns}
  \note{
    \textbf{Recadrage}
    \begin{itemize}
      \item Paysage connu depuis leçon 1, mais présenté ici comme un \emph{pari}.
      \item Garantie \emph{mathématique} vs. statistique (tests).
    \end{itemize}
    \textbf{Ordres de grandeur du goulot humain}
    \begin{itemize}
      \item Preuve math $\approx 10$--$20\times$ le coût de la rédaction informelle.
      \item \emph{Calculemus} de Leibniz --- à reprendre en clôture.
    \end{itemize}
    \textbf{Question qui structure l'acte I}
    \begin{itemize}
      \item Et si l'IA venait lever ce goulot ?
    \end{itemize}
  }
\end{frame}

% ─────────────────────────────────────────────────────────────
% FRAME 3 — L'IA GÉNÈRE, LE NOYAU VÉRIFIE (TIKZ)
% ─────────────────────────────────────────────────────────────
\begin{frame}{La nouvelle division du travail : l'IA génère, le noyau vérifie}
  \begin{center}
    \begin{tikzpicture}[
        >=Stealth,
        box/.style={draw, rounded corners=4pt, align=center, minimum height=1.1cm, font=\small},
        arr/.style={->, thick},
        lbl/.style={font=\scriptsize, text=gray}
      ]
      \node[box, fill=blue!15, text width=3cm] (human) at (-5.5, 0) {Humain\\{\scriptsize spécification, intuition}};
      \node[box, fill=orange!15, text width=3cm] (llm) at (0, 0) {LLM / IA\\{\scriptsize génère des tentatives de preuve}};
      \node[box, fill=green!15, text width=3cm] (kernel) at (5.5, 0) {Noyau Rocq/Lean\\{\scriptsize vérifie $\lambda$-terme}};

      \draw[arr] (human) -- node[above, lbl] {énoncé} node[below, lbl] {\texttt{Theorem T : \ldots}} (llm);
      \draw[arr] (llm) -- node[above, lbl] {tactiques / terme} (kernel);
      \draw[arr, dashed] (kernel.south) to[bend left=30] node[below, lbl] {échec / feedback} (llm.south);
      \draw[arr, dashed] (kernel.north) to[bend right=40] node[above=4pt, lbl] {certificat vérifié} (human.north);
    \end{tikzpicture}
  \end{center}
  \begin{itemize}
    \item \textbf{Autoformalisation} : informel $\to$ Rocq/Lean~\cite{buzzardMathematicalReasoningComputer2025}
    \item \textbf{Tactic search} : AlphaProof~\cite{alphaproof2024}
    \item \textbf{L'IA propose, le noyau dispose}
  \end{itemize}
  \note{
    \textbf{Asymétrie à souligner}
    \begin{itemize}
      \item LLM faillible : hallucinations, lemmes inventés $\to$ sans conséquence.
      \item Noyau rejette tout $\lambda$-terme mal typé.
      \item Inédit parmi les applications d'IA : pas besoin de faire confiance au modèle.
    \end{itemize}
    \textbf{Pont leçon 5}
    \begin{itemize}
      \item \texttt{Fixpoint} bien typé $\Rightarrow$ terminaison garantie, peu importe l'auteur.
    \end{itemize}
    \textbf{Nuance à glisser}
    \begin{itemize}
      \item Division du travail robuste \emph{en apparence} --- acte II reviendra dessus.
    \end{itemize}
  }
\end{frame}

% ─────────────────────────────────────────────────────────────
% FRAME 4 — ÉTAT DE L'ART : MATHÉMATIQUES
% ─────────────────────────────────────────────────────────────
\begin{frame}{Mathématiques formalisées : l'appétit de l'IA}
  \begin{columns}
    \begin{column}{0.55\textwidth}
      \textbf{Avigad}~\cite{avigadMathematicsFormalTurn2023} — \emph{tournant formel}
      \smallskip

      \textbf{Buzzard}~\cite{buzzardMathematicalReasoningComputer2025,leanlearningseminarProfessorKevinBuzzard2025} —
      après le calcul, le \emph{raisonnement} ?
      \medskip

      \textbf{Jalons :}
      \begin{itemize}
        \item Liquid Tensor~\cite{scholze2021liquid} / Mathlib
        \item AlphaProof, IMO 2024~\cite{alphaproof2024}
        \item \textbf{2026} : problèmes d'Erd\H{o}s~\cite{erdosproblems} résolus\\
          par Aristotle~\cite{aristotle2026erdos728}, Aletheia~\cite{deepmind2026aletheia}
      \end{itemize}
    \end{column}
    \begin{column}{0.42\textwidth}
      \begin{alertblock}{Severini~\cite{severini2026infrastructure}}
        \small
        $\approx 0{,}05\%$ des maths formalisées.\\
        L'IA a faim de corpus vérifié : la formalisation devient \emph{application phare}.
      \end{alertblock}
      \medskip
      \small\emph{Massot~\cite{massotWhyFormalizeMathematics} : outil, pas substitut.}
    \end{column}
  \end{columns}
  \note{
    \textbf{Seuil qualitatif fin 2025 -- début 2026}
    \begin{itemize}
      \item Passage \emph{maths de compétition} (IMO) $\to$ \emph{maths de recherche}.
      \item Janv.\ 2026 : problèmes \#397, \#728, \#729 résolus en une semaine (GPT-5.2 Pro + Aristotle).
      \item Févr.\ 2026 : Aletheia (Gemini 3 Deep Think) résout 4 problèmes dont \#1051.
      \item $\approx 100$ énoncés basculés en \og résolus~\fg{} depuis octobre 2025.
    \end{itemize}
    \textbf{Severini --- inversion du moteur}
    \begin{itemize}
      \item Moteur plus seulement scientifique : aussi \emph{économique/infrastructurel}.
      \item Labos d'IA ont un intérêt direct à financer la formalisation.
    \end{itemize}
    \textbf{Nuance Massot}
    \begin{itemize}
      \item Formalisation = outil, pas substitut.
      \item Pluralité de styles (informel / semi-formel / mécanisé) qui coexistent.
    \end{itemize}
  }
\end{frame}

% ─────────────────────────────────────────────────────────────
% FRAME 5 — EXTENSION : SÛRETÉ DE L'IA
% ─────────────────────────────────────────────────────────────
\begin{frame}{Au-delà des maths : sûreté de l'IA par la preuve}
  \begin{block}{Tegmark \& Omohundro~\cite{tegmark2023provably}}
    \emph{Provably Safe Systems} — seule voie vers une AGI contrôlable.
  \end{block}
  \pause
  \medskip
  \begin{enumerate}
    \item Alignement (RLHF, monitoring) $\to$ garanties \emph{statistiques}
    \item Preuve formelle $\to$ garantie \emph{mathématique}
    \item L'IA rend l'approche \emph{faisable} à grande échelle
  \end{enumerate}
  \pause
  \medskip
  \begin{alertblock}{La promesse ultime}
    Vérification $+$ IA $=$ sûreté des systèmes critiques, \emph{y compris} des IA.
  \end{alertblock}
  \note{
    \textbf{Reconstruire l'argument à voix haute}
    \begin{itemize}
      \item RLHF, monitoring : garanties \emph{statistiques}.
      \item À grande échelle, \og 1 cas sur 1000~\fg{} devient inacceptable.
      \item Seule la preuve formelle fournit l'assurance requise.
      \item Ce qui rend l'approche faisable : c'est l'IA qui rédige les preuves.
    \end{itemize}
    \textbf{Tension à annoncer}
    \begin{itemize}
      \item Prouver \og pas de stack overflow~\fg{} $\neq$ prouver \og ne pas nuire~\fg.
      \item Le problème de la \emph{spécification} reste entier (traité à l'acte II).
    \end{itemize}
    \textbf{Transition vers acte II}
    \begin{itemize}
      \item Et si le cadre technique lui-même était fragile ?
    \end{itemize}
  }
\end{frame}

% ─────────────────────────────────────────────────────────────
% ACTE II — LE RETOURNEMENT
% ─────────────────────────────────────────────────────────────
\section{Acte II — Le retournement}

% ─────────────────────────────────────────────────────────────
% FRAME 6 — DE BRUIJN : QUI VÉRIFIE LE NOYAU ?
% ─────────────────────────────────────────────────────────────
\begin{frame}{Qui vérifie le noyau ?}
  \small\emph{Si on fonde la sûreté sur la preuve, tout repose sur le vérificateur —
  commençons par lui.}
  \medskip

  \begin{block}{Critère de de Bruijn~\cite{debruijn1980automath}}
    Noyau \emph{minimal} : vérificateur de $\lambda$-termes, auditable par un humain.
  \end{block}
  \medskip

  \begin{itemize}
    \item Pas de confiance dans les tactiques
    \item Confiance dans le \textbf{noyau} ($\sim$ quelques milliers de lignes)
    \item Noyau correct $\Rightarrow$ toute preuve acceptée est correcte
  \end{itemize}
  \pause
  \medskip

  \begin{alertblock}{Régression}
    Noyau correct ? \quad Compilateur du noyau ? \quad Matériel ?
  \end{alertblock}
  \note{
    \textbf{Transition courte --- ne pas s'attarder}
    \begin{itemize}
      \item Si la sûreté repose sur la preuve : examiner le vérificateur.
    \end{itemize}
    \textbf{Point d'architecture à expliciter}
    \begin{itemize}
      \item Tactiques arbitrairement complexes $\to$ pas grave (produisent des termes).
      \item Toute la confiance est concentrée dans le noyau.
    \end{itemize}
    \textbf{Régression à tracer au tableau}
    \begin{itemize}
      \item Noyau $\to$ compilateur du noyau $\to$ matériel.
      \item Thématisée par Thompson (1984) --- slide suivante.
    \end{itemize}
  }
\end{frame}

% ─────────────────────────────────────────────────────────────
% FRAME 7 — THOMPSON 1984
% ─────────────────────────────────────────────────────────────
\begin{frame}{Thompson 1984 — \emph{Reflections on Trusting Trust}}
  \begin{columns}
    \begin{column}{0.55\textwidth}
      Ken Thompson, Turing Award~\cite{thompson1984trusting}.
      \medskip

      \textbf{Compilateur C malveillant :}
      \begin{enumerate}
        \item porte dérobée injectée dans \texttt{login}
        \item s'auto-reproduit à la compilation d'un compilateur
      \end{enumerate}
      \medskip

      Source \emph{propre} $\to$ binaire \emph{compromis}.
    \end{column}
    \begin{column}{0.42\textwidth}
      \begin{block}{Thèse}
        \emph{« You can't trust code that you did not totally create yourself. »}
      \end{block}
      \medskip
      \begin{alertblock}{Conséquence}
        La vérification \emph{déplace} la confiance — elle ne l'élimine pas.
      \end{alertblock}
    \end{column}
  \end{columns}
  \note{
    \textbf{Portée de l'argument}
    \begin{itemize}
      \item Pas un bug particulier : \emph{structure générale} de la confiance en informatique.
      \item Origine de la notion de \emph{trusted computing base}.
    \end{itemize}
    \textbf{Justification a posteriori du critère de de Bruijn}
    \begin{itemize}
      \item Minimiser le noyau = minimiser la surface de confiance.
      \item Mais cette surface n'est \emph{jamais nulle} --- un petit noyau s'exécute toujours quelque part de non prouvé.
    \end{itemize}
    \textbf{Préparer la chute}
    \begin{itemize}
      \item 2026 donne des exemples concrets à chaque niveau qui craque.
    \end{itemize}
  }
\end{frame}

% ─────────────────────────────────────────────────────────────
% FRAME 8 — LE RETOURNEMENT 2026 (FULL SLIDE)
% ─────────────────────────────────────────────────────────────
\begin{frame}{Le retournement 2026 : l'IA casse le cadre}
  \begin{columns}[T]
    \begin{column}{0.32\textwidth}
      \begin{block}{Mythos~\cite{anthropic2026mythos}}
        \scriptsize
        \textbf{Anthropic, avril 2026}
        \smallskip

        $\bullet$ 0-days dans OpenBSD, Linux, FreeBSD\\
        $\bullet$ Milliers de vulnérabilités\\
        $\bullet$ Source reconstruite depuis binaires
      \end{block}
      \scriptsize\emph{L'IA comme attaquant de classe nouvelle.}
    \end{column}
    \begin{column}{0.32\textwidth}
      \begin{alertblock}{\texttt{False} dans Rocq~\cite{rocqfalse2026,sterin2026falsehood}}
        \scriptsize
        \textbf{Mars 2026}
        \smallskip

        $\bullet$ Sept preuves de \texttt{False} via bugs indépendants\\
        $\bullet$ Guard checker, modules\\
        $\bullet$ IA $+$ humain, avec l'équipe Rocq
      \end{alertblock}
      \scriptsize\emph{Le noyau minimal n'est pas infaillible.}
    \end{column}
    \begin{column}{0.32\textwidth}
      \begin{exampleblock}{\emph{Watchers}~\cite{gopinathan2026watchers}}
        \scriptsize
        \textbf{\texttt{lean-zip}, avril 2026}
        \smallskip

        $\bullet$ Décodeur ZIP prouvé en Lean 4\\
        $\bullet$ Parseur non vérifié : overflow\\
        $\bullet$ Un 2\textsuperscript{e} bug hors-spec\\
        $\bullet$ Trouvés par Claude Code
      \end{exampleblock}
      \scriptsize\emph{Preuve OK, runtime cassé, spec incomplète.}
    \end{column}
  \end{columns}
  \medskip
  \begin{center}
    \textbf{Trois niveaux, trois fractures --- toutes mises au jour par l'IA.}
  \end{center}
  \note{
    \textbf{\underline{Slide centrale --- prendre le temps.}}
    \smallskip

    \textbf{Mythos}
    \begin{itemize}
      \item Techniques : RCE, ROP, JIT ; bugs retrouvés 27 ans après leur introduction.
      \item \emph{Retourne Tegmark} : l'IA qui vérifie est aussi celle qui casse --- course, non victoire.
    \end{itemize}
    \textbf{\texttt{False} dans Rocq}
    \begin{itemize}
      \item Bugs exploités : \emph{guard checker} + gestion des modules.
      \item Collaboration humains $+$ Claude Opus 4.6 (Stérin), concertée avec l'équipe Rocq.
      \item Distinction clé : \emph{auditabilité} $\neq$ \emph{correction}.
    \end{itemize}
    \textbf{Watchers}
    \begin{itemize}
      \item Décodeur prouvé ; le code \emph{appelant} (\texttt{lean\_alloc\_sarray}) ne l'est pas $\to$ heap overflow.
      \item Second bug : comportement exploité \emph{non couvert} par la spécification --- ce n'est donc pas seulement un problème de runtime.
      \item Découverte quasi autonome par \emph{Claude Code} (rôle humain : protocole + rapport).
      \item Thompson 1984 littéralement appliqué à la vérification formelle.
    \end{itemize}
    \textbf{Point commun à souligner à la fin}
    \begin{itemize}
      \item C'est \emph{l'IA} qui met au jour les trois fractures.
    \end{itemize}
  }
\end{frame}

% ─────────────────────────────────────────────────────────────
% FRAME 9 — LES 3 NIVEAUX DE CONFIANCE
% ─────────────────────────────────────────────────────────────
\begin{frame}{Trois niveaux de confiance, ré-examinés}
  \begin{center}
    \begin{tikzpicture}[
        >=Stealth,
        layer/.style={draw, rounded corners=4pt, text width=7cm, align=left,
                      minimum height=1cm, font=\small},
        broken/.style={text=red!70!black, font=\scriptsize\itshape}
      ]
      \node[layer, fill=blue!10] (spec) at (0, 2.4)
        {\textbf{1. Spécification}};
      \node[layer, fill=orange!10] (kernel) at (0, 1.0)
        {\textbf{2. Noyau}};
      \node[layer, fill=green!10] (runtime) at (0, -0.4)
        {\textbf{3. Runtime / compilateur / matériel}};

      \node[broken, right=0.2cm of spec.east, text width=4cm] {Formaliser les valeurs ?};
      \node[broken, right=0.2cm of kernel.east, text width=4cm] {\texttt{False} (Opus 4.6)~\cite{rocqfalse2026}};
      \node[broken, right=0.2cm of runtime.east, text width=4cm] {Overflow Lean~\cite{gopinathan2026watchers}};
    \end{tikzpicture}
  \end{center}
  \medskip
  \begin{alertblock}{À retenir}
    Pas de garantie absolue. Garanties \emph{relatives} à des hypothèses explicites.
  \end{alertblock}
  \note{
    \textbf{Recadrage, pas réfutation}
    \begin{itemize}
      \item Vérification formelle reste l'outil le plus puissant pour les garanties logicielles.
      \item Ses garanties sont \emph{conditionnelles} à des hypothèses explicites.
    \end{itemize}
    \textbf{Analogie à développer}
    \begin{itemize}
      \item Prouver en théorie des ensembles = prouver \emph{si} ZFC cohérent.
      \item Prouver en Rocq = prouver \emph{si} noyau correct $\wedge$ runtime fidèle $\wedge$ matériel sain.
    \end{itemize}
    \textbf{Formule clé à noter au tableau}
    \begin{itemize}
      \item Pas \og peut-on faire confiance ?~\fg.
      \item Mais \og à quoi fait-on confiance, et est-ce justifié ?~\fg.
    \end{itemize}
  }
\end{frame}

% ─────────────────────────────────────────────────────────────
% ACTE III — QUESTIONS OUVERTES
% ─────────────────────────────────────────────────────────────
\section{Acte III — Questions ouvertes}

% ─────────────────────────────────────────────────────────────
% FRAME 10 — QUESTIONS THINK-PAIR-SHARE
% ─────────────────────────────────────────────────────────────
% \begin{frame}{Discussion — Think/Pair/Share}
%   \small
%   \textbf{Format :} 2 min seul, 3 min en binôme, mise en commun.
%   \medskip

%   \begin{enumerate}
%     \item \textbf{IA : adversaire ou garante ?}\\
%       Red team systématique, ou perte de confiance ?
%     \smallskip

%     \item \textbf{Formaliser Spinoza, Gödel}~\cite{pommeret2025spinoza,benzmuller2013goedel,kirchner2017zalta}\\
%       Traduction fidèle, ou ré-interprétation ?
%     \smallskip

%     \item \textbf{Comprendre vs. vérifier}\\
%       Preuve de 10\,000 lignes jamais lue — qui la comprend ?
%     \smallskip

%     \item \textbf{Sûreté prouvable~\cite{tegmark2023provably} : voie royale ou illusion ?}\\
%       Peut-on prouver des valeurs non-formelles ?
%   \end{enumerate}
%   \note{
%     \textbf{Gestion du temps}
%     \begin{itemize}
%       \item 2 min seul / 3 min binôme / mise en commun.
%       \item 2 questions max sur 4 --- choisir selon l'énergie de la salle.
%     \end{itemize}
%     \textbf{Profil des questions}
%     \begin{itemize}
%       \item Q1 : concret (sécurité, red team, course offensive/défensive).
%       \item Q2 : philosophique (traduction fidèle vs. ré-interprétation).
%       \item Q3 : épistémologique --- Gonthier (4 couleurs) et Hales (Kepler) : preuves jamais lues intégralement, seulement vérifiées. \emph{Est-ce faire des maths ?}
%       \item Q4 : IA safety / méta-éthique (peut-on formaliser des valeurs ?).
%     \end{itemize}
%   }
% \end{frame}

% ─────────────────────────────────────────────────────────────
% FRAME 11 — BILAN DU COURS
% ─────────────────────────────────────────────────────────────
\begin{frame}{Bilan — six semaines}
  \begin{columns}[T]
    \begin{column}{0.48\textwidth}
      \textbf{Savoir-faire :}
      \begin{itemize}
        \item Habiter un type
        \item Déduction naturelle
        \item Rocq : \texttt{Inductive}, \texttt{Fixpoint}, \texttt{induction}
      \end{itemize}
    \end{column}
    \begin{column}{0.48\textwidth}
      \textbf{Philosophiquement :}
      \begin{itemize}
        \item \textbf{Curry-Howard} : preuves $=$ programmes
        \item \textbf{Constructivisme}
        \item \textbf{Vérificationnisme}
        \item \textbf{Mécanisation} du travail épistémique
        \item \textbf{Confiance} déplacée, pas supprimée
      \end{itemize}
      \medskip
      \begin{block}{Pour la suite}
        \small
        DM — rev $\circ$ rev $=$ id — 25 avril.\\
        \emph{Calculemus}, mais lucidement.
      \end{block}
    \end{column}
  \end{columns}
  \note{
    \textbf{Insister}
    \begin{itemize}
      \item Double niveau (technique $+$ philosophique) tenu sur 6 semaines.
      \item \og Confiance déplacée, pas supprimée~\fg{} --- thèse centrale du cours.
    \end{itemize}
    \textbf{Cadre historique}
    \begin{itemize}
      \item Tournant formel s'accélère ; questions épistémologiques loin d'être résolues.
      \item Les philosophes ont leur mot à dire dans ce débat.
    \end{itemize}
    \textbf{Clore}
    \begin{itemize}
      \item DM : \texttt{rev (rev l) = l}, rendu le 25 avril.
      \item \emph{Calculemus, mais lucidement.}
      \item Remerciements.
    \end{itemize}
  }
\end{frame}

% ─────────────────────────────────────────────────────────────
% FRAME 12 — BIBLIOGRAPHIE
% ─────────────────────────────────────────────────────────────
\begin{frame}[allowframebreaks]{Bibliographie}
  \tiny
  \printbibliography[heading=none]
\end{frame}

\end{document}
